\documentclass[a4paper,12pt]{article}

\usepackage[T2A]{fontenc}
\usepackage[utf8]{inputenc}
\usepackage[russian]{babel}

\usepackage{amsmath, amssymb, amsthm, mathtools}
\usepackage{helvet}
\renewcommand{\familydefault}{\sfdefault}
\usepackage{icomma}
\usepackage{geometry}
\usepackage{microtype}
\usepackage{tikz}
\usepackage{tkz-euclide}
\usepackage{pgfplots}
\pgfplotsset{compat=1.18}
\usepackage{tabularx, booktabs, longtable}
\usepackage{enumitem}
\usepackage{hyperref}
\usepackage{parskip}
\usepackage{fancyhdr}
\usepackage{setspace}
\usepackage{xcolor}

\definecolor{accent}{HTML}{008080}
\definecolor{darkaccent}{HTML}{004D4D}
\definecolor{textgray}{HTML}{333333}
\definecolor{softgray}{HTML}{F3F6F6}

\geometry{left=18mm,right=18mm,top=18mm,bottom=20mm}
\onehalfspacing
\setlength{\parindent}{0pt}
\setlist[itemize]{leftmargin=1.2cm,itemsep=2pt,topsep=3pt}
\setlist[enumerate]{leftmargin=1.2cm,itemsep=3pt,topsep=3pt}

\hypersetup{
    colorlinks=true,
    linkcolor=darkaccent,
    urlcolor=darkaccent
}

\pagestyle{fancy}
\fancyhf{}
\fancyhead[L]{\textcolor{textgray}{Показательные неравенства}}
\fancyhead[R]{\textcolor{textgray}{Чек-лист}}
\fancyfoot[C]{\thepage}
\renewcommand{\headrulewidth}{0.4pt}
\renewcommand{\headrule}{\hbox to\headwidth{\color{accent}\leaders\hrule height \headrulewidth\hfill}}

\newcommand{\sectionline}{\vspace{2mm}{\color{accent}\hrule height 0.6pt}\vspace{2mm}}
\newcommand{\important}[1]{\textcolor{darkaccent}{\textbf{#1}}}
\newcommand{\R}{\mathbb{R}}

\begin{document}

\begin{titlepage}
\thispagestyle{empty}
\centering

\vspace*{28mm}

{\Huge\bfseries\textcolor{darkaccent}{Чек-лист}\par}
\vspace{5mm}
{\LARGE\bfseries Показательные неравенства\par}

\vspace{18mm}

{\large Лёвин Артём Александрович эксклюзивно для Николь\par}

\vspace{8mm}

{\large \today\par}

\vfill

\begin{tikzpicture}[x=1cm,y=1cm]
\draw[accent,line width=1.1pt]
    (-6,0) .. controls (-3.5,1.1) and (-1.5,-1.1) .. (0,0)
           .. controls (1.5,1.1) and (3.5,-1.1) .. (6,0);
\draw[darkaccent,line width=0.5pt]
    (-4.8,-0.45) .. controls (-2.2,0.45) and (2.2,-0.45) .. (4.8,0.45);
\end{tikzpicture}

\vspace*{12mm}
\end{titlepage}

\section*{1. Главная идея показательных неравенств}
\sectionline

Показательное неравенство содержит переменную в показателе степени:
\[
a^{f(x)} \ge a^{g(x)},\qquad a>0,\quad a\ne 1.
\]

\important{Ключевая цель} --- привести левую и правую части к одинаковому основанию или сделать замену.

\begin{center}
\renewcommand{\arraystretch}{1.35}
\begin{tabularx}{\linewidth}{@{}p{0.31\linewidth}X@{}}
\toprule
\textbf{Ситуация} & \textbf{Что делать} \\
\midrule
Слева и справа по одной степени & Привести к одному основанию и сравнить показатели. \\
Есть суммы, разности, дроби со степенями & Ищем повторяющийся элемент и делаем замену. \\
После замены появилась дробь & Решаем рациональное неравенство методом интервалов. \\
После обратной замены нужно сравнить с числом & При необходимости используем логарифмы. \\
\bottomrule
\end{tabularx}
\end{center}

\section*{2. Сравнение степеней с одинаковым основанием}
\sectionline

Пусть \(a>0\), \(a\ne 1\). Тогда:

\[
a^{f(x)} \ge a^{g(x)}.
\]

\begin{itemize}
    \item Если \(a>1\), то функция \(a^x\) возрастает, знак неравенства сохраняется:
    \[
    a^{f(x)} \ge a^{g(x)} \Longleftrightarrow f(x)\ge g(x).
    \]
    \item Если \(0<a<1\), то функция \(a^x\) убывает, знак неравенства меняется:
    \[
    a^{f(x)} \ge a^{g(x)} \Longleftrightarrow f(x)\le g(x).
    \]
\end{itemize}

\important{Нельзя просто ``откинуть основание'' без проверки, больше оно единицы или меньше единицы.}

\subsection*{Мини-пример}

\[
\left(\frac{\pi}{3}\right)^{|x^2-x|-3}\ge 1.
\]

Так как
\[
1=\left(\frac{\pi}{3}\right)^0,\qquad \frac{\pi}{3}>1,
\]
то знак сохраняется:
\[
|x^2-x|-3\ge 0.
\]

\section*{3. Правила для модуля}
\sectionline

Перед раскрытием модуля его нужно изолировать:
\[
|f(x)| \ \text{слева},\qquad g(x)\ \text{справа}.
\]

\subsection*{Случай 1. Модуль больше}

\[
|f(x)|\ge g(x)
\]

при \(g(x)\ge 0\) удобно раскрывать как совокупность:
\[
\left[
\begin{array}{l}
f(x)\ge g(x),\\
f(x)\le -g(x).
\end{array}
\right.
\]

Если \(g(x)<0\), то неравенство \(|f(x)|\ge g(x)\) выполняется при всех допустимых \(x\), потому что \(|f(x)|\ge 0\).

\subsection*{Случай 2. Модуль меньше}

\[
|f(x)|\le g(x)
\]

эквивалентен системе:
\[
\left\{
\begin{array}{l}
f(x)\le g(x),\\
f(x)\ge -g(x).
\end{array}
\right.
\]

Здесь обязательно нужно, чтобы правая часть была неотрицательной. Если \(g(x)<0\), решений нет.

\section*{4. Система и совокупность}
\sectionline

\begin{center}
\renewcommand{\arraystretch}{1.35}
\begin{tabularx}{\linewidth}{@{}p{0.22\linewidth}X@{}}
\toprule
\textbf{Запись} & \textbf{Смысл} \\
\midrule
Система \(\left\{\begin{array}{l} A\\ B \end{array}\right.\) &
Должны выполняться оба условия одновременно. В ответ идёт пересечение решений. \\
\addlinespace
Совокупность \(\left[\begin{array}{l} A\\ B \end{array}\right.\) &
Достаточно выполнения хотя бы одного условия. В ответ идёт объединение решений. \\
\bottomrule
\end{tabularx}
\end{center}

\important{Практический ориентир:}
\begin{itemize}
    \item \(|f(x)|\ge c\) обычно даёт совокупность;
    \item \(|f(x)|\le c\) обычно даёт систему;
    \item отрезок вида \(a\le t\le b\) при обратной замене записывается как система.
\end{itemize}

\section*{5. Метод интервалов для квадратных неравенств}
\sectionline

Квадратное неравенство обычно приводят к виду:
\[
ax^2+bx+c \ge 0
\quad \text{или} \quad
ax^2+bx+c \le 0.
\]

Алгоритм:
\begin{enumerate}
    \item Заменить неравенство уравнением:
    \[
    ax^2+bx+c=0.
    \]
    \item Найти корни.
    \item Отметить корни на числовой прямой.
    \item Закрасить точки, если неравенство нестрогое: \(\le\), \(\ge\).
    \item Расставить знаки на промежутках.
    \item Выбрать промежутки с нужным знаком.
\end{enumerate}

\begin{center}
\begin{tikzpicture}[scale=1]
\draw[->,line width=0.8pt] (-4,0)--(4,0);
\foreach \x/\lab in {-2/{$x_1$},2/{$x_2$}}{
    \filldraw[accent] (\x,0) circle (2pt);
    \node[below] at (\x,-0.05) {\lab};
}
\node[above] at (-3,0.1) {$+$};
\node[above] at (0,0.1) {$-$};
\node[above] at (3,0.1) {$+$};
\draw[accent,line width=1.2pt] (-4,0.18)--(-2,0.18);
\draw[accent,line width=1.2pt] (2,0.18)--(4,0.18);
\end{tikzpicture}

\small Для параболы с положительным старшим коэффициентом знак \(+\) обычно по краям.
\end{center}

\section*{6. Приведение к общему основанию}
\sectionline

Используем свойства степеней:
\[
a^m\cdot a^n=a^{m+n},\qquad
\frac{a^m}{a^n}=a^{m-n},\qquad
(a^m)^n=a^{mn}.
\]

\subsection*{Типовые преобразования}

\[
0{,}5=\frac12=2^{-1},\qquad
8=2^3,\qquad
9=3^2,\qquad
49=7^2.
\]

\[
0{,}5^{\,x}\cdot 2^{x^2}
=
(2^{-1})^x\cdot 2^{x^2}
=
2^{-x}\cdot 2^{x^2}
=
2^{x^2-x}.
\]

\important{Полезное правило:} лучше приводить к основанию, которое больше \(1\). Тогда при сравнении показателей знак неравенства не придётся менять.

\section*{7. Замена в показательных неравенствах}
\sectionline

Замена нужна, когда:
\begin{itemize}
    \item одинаковое основание встречается несколько раз;
    \item степени отличаются на константу;
    \item есть выражения вида \(3^x\), \(9^x\), \(27^x\);
    \item есть сумма, разность или дробь, из-за которых нельзя просто сравнить показатели.
\end{itemize}

\subsection*{Главное ограничение}

Если
\[
t=a^{f(x)},\qquad a>0,\quad a\ne 1,
\]
то всегда
\[
t>0.
\]

Это ограничение обязательно учитывается при решении неравенства относительно \(t\).

\subsection*{Пример замены}

\[
9^x=(3^2)^x=(3^x)^2.
\]

Если обозначить
\[
t=3^x,\qquad t>0,
\]
то
\[
9^x=t^2.
\]

\section*{8. Обратная замена и логарифмы}
\sectionline

После решения неравенства относительно \(t\) нужно вернуться к \(x\).

Если
\[
t=3^x
\]
и получилось, например,
\[
\frac12\le t<1,
\]
то нужно записать:
\[
\frac12\le 3^x<1.
\]

Для приведения чисел к основанию \(3\) используется основное логарифмическое тождество:
\[
b=a^{\log_a b},\qquad a>0,\quad a\ne 1,\quad b>0.
\]

Например:
\[
\frac12=3^{\log_3\frac12},\qquad
\sqrt2=3^{\log_3\sqrt2},\qquad
1=3^0.
\]

Так как \(3>1\), знак сохраняется:
\[
3^x\ge 3^{\log_3\frac12}
\Longleftrightarrow
x\ge \log_3\frac12.
\]

\section*{9. Рациональное неравенство после замены}
\sectionline

Иногда после замены получается дробное неравенство:
\[
\frac{7}{t^2-2}-\frac{2}{t-1}\ge 0.
\]

Алгоритм:
\begin{enumerate}
    \item Привести к общему знаменателю.
    \item Найти нули числителя.
    \item Найти нули знаменателя.
    \item Нули знаменателя обязательно выколоть.
    \item Учесть ограничение замены, например \(t>0\).
    \item Решить методом интервалов.
    \item Перейти обратно от \(t\) к \(x\).
\end{enumerate}

\important{Нельзя сокращать множители в неравенстве без контроля знака и ОДЗ.} Безопаснее решать через общий знаменатель и метод интервалов.

\section*{10. Разбор ключевого примера}
\sectionline

Решим неравенство:
\[
\left(\frac{\pi}{3}\right)^{|x^2-x|-3}\ge 1.
\]

\textbf{Шаг 1. Приводим правую часть к тому же основанию.}
\[
1=\left(\frac{\pi}{3}\right)^0.
\]

\textbf{Шаг 2. Проверяем основание.}
\[
\frac{\pi}{3}\approx 1{,}047>1.
\]
Знак неравенства сохраняется:
\[
|x^2-x|-3\ge 0.
\]

\textbf{Шаг 3. Изолируем модуль.}
\[
|x^2-x|\ge 3.
\]

\textbf{Шаг 4. Раскрываем модуль через совокупность.}
\[
\left[
\begin{array}{l}
x^2-x\ge 3,\\
x^2-x\le -3.
\end{array}
\right.
\]

\textbf{Шаг 5. Решаем первое неравенство.}
\[
x^2-x-3\ge 0.
\]
Корни:
\[
x_{1,2}=\frac{1\pm\sqrt{13}}{2}.
\]
Так как ветви параболы направлены вверх:
\[
x\in
\left(-\infty,\frac{1-\sqrt{13}}{2}\right]
\cup
\left[\frac{1+\sqrt{13}}{2},+\infty\right).
\]

\textbf{Шаг 6. Решаем второе неравенство.}
\[
x^2-x+3\le 0.
\]
Дискриминант:
\[
D=1-12=-11<0.
\]
Парабола всегда положительна, поэтому решений нет:
\[
\emptyset.
\]

\textbf{Ответ:}
\[
x\in
\left(-\infty,\frac{1-\sqrt{13}}{2}\right]
\cup
\left[\frac{1+\sqrt{13}}{2},+\infty\right).
\]

\section*{11. Типичные ошибки}
\sectionline

\begin{enumerate}
    \item Не проверить основание степени и не поменять знак при \(0<a<1\).
    \item Записать \(1=a^1\) вместо правильного \(1=a^0\).
    \item Раскрыть \(|f(x)|\ge c\) как систему вместо совокупности.
    \item Раскрыть \(|f(x)|\le c\) как совокупность вместо системы.
    \item Забыть, что при замене \(t=a^{f(x)}\) всегда \(t>0\).
    \item Отметить на оси \(t\) отрицательные значения, хотя \(t>0\).
    \item Не выколоть корни знаменателя в рациональном неравенстве.
    \item Сократить выражение в неравенстве и потерять ОДЗ.
    \item Перепутать объединение и пересечение промежутков.
    \item После решения относительно \(t\) забыть сделать обратную замену.
\end{enumerate}

\section*{12. Краткий алгоритм решения}
\sectionline

\begin{enumerate}
    \item Определите тип неравенства:
    \begin{itemize}
        \item можно привести к одному основанию;
        \item нужна замена;
        \item после замены будет рациональное неравенство.
    \end{itemize}
    \item Если основания одинаковые, проверьте:
    \[
    a>1 \quad \text{или} \quad 0<a<1.
    \]
    \item Если есть модуль, сначала изолируйте его.
    \item Если делаете замену, обязательно запишите ограничение \(t>0\).
    \item Решите полученное алгебраическое неравенство.
    \item Переведите ответ обратно к переменной \(x\).
    \item Проверьте ОДЗ, выколотые точки и смысл ответа.
\end{enumerate}

\newpage

\section*{Тренировочный блок}
\sectionline

Решите неравенства. Упражнения расположены по нарастающей сложности.

\subsection*{Средний уровень}

\begin{enumerate}
    \item
    \[
    \left(\frac{\pi}{3}\right)^{|x^2+2x|-4}\ge 1.
    \]

    \item
    \[
    2^{x^2-3x}\le 8.
    \]

    \item
    \[
    \left(\frac12\right)^{x^2-5x+4}\ge \left(\frac12\right)^2.
    \]
\end{enumerate}

\subsection*{Выше среднего}

\begin{enumerate}[resume]
    \item
    \[
    0{,}5^{\,2x}\cdot 4^{x-1}>8^{x-2}.
    \]

    \item
    \[
    9^x-4\cdot 3^x+3\le 0.
    \]

    \item
    \[
    4^x-5\cdot 2^x+4>0.
    \]

    \item
    \[
    \frac{6}{3^x-1}\ge 2.
    \]

    \item
    \[
    \frac{5}{9^x-4}-\frac{1}{3^x-2}\le 0.
    \]

    \item
    \[
    7^{x^2+2x-3}-8\cdot 7^{x^2+2x-4}\ge 0.
    \]
\end{enumerate}

\subsection*{Сложный уровень}

\begin{enumerate}[resume]
    \item
    \[
    98-7^{x^2+5x-48}\ge 49^{x^2+5x-49}.
    \]
\end{enumerate}

\end{document}